% Non-Associative State Space Models: Octonion Dynamics for Path-Dependent Sequence Modeling
% NeurIPS 2026 submission
%
% Compile with: pdflatex paper_a_ossm.tex && bibtex paper_a_ossm && pdflatex paper_a_ossm.tex && pdflatex paper_a_ossm.tex

\documentclass{article}

% NeurIPS 2026 style
\usepackage[final]{neurips_2026}

\usepackage[utf8]{inputenc}
\usepackage[T1]{fontenc}
\usepackage{hyperref}
\usepackage{url}
\usepackage{booktabs}
\usepackage{amsfonts}
\usepackage{amsmath}
\usepackage{amssymb}
\usepackage{nicefrac}
\usepackage{microtype}
\usepackage{graphicx}
\usepackage{xcolor}
\usepackage{mathtools}
\usepackage{algorithm}
\usepackage{algorithmic}
\usepackage{amsthm}

\newtheorem{theorem}{Theorem}
\usepackage{tikz}
\usepackage{multirow}
\usepackage{caption}

% Custom commands
\newcommand{\Oct}{\mathbb{O}}
\newcommand{\R}{\mathbb{R}}
\newcommand{\C}{\mathbb{C}}
\newcommand{\HH}{\mathbb{H}}
\newcommand{\Sed}{\mathbb{S}}
\newcommand{\ossm}{\text{O-SSM}}
\newcommand{\assoc}[3]{[#1,#2,#3]}
\newcommand{\norm}[1]{\lvert #1 \rvert}
\newcommand{\abs}[1]{\lvert #1 \rvert}
\newcommand{\todo}[1]{{\color{red}\textbf{[TODO: #1]}}}

\title{Non-Associative State Space Models: Octonion Dynamics\\for Path-Dependent Sequence Modeling}

\author{%
  Demetrios Chiuratto Agourakis\textsuperscript{1}\quad
  Dionisio Chiuratto Agourakis\textsuperscript{2}\\[4pt]
  \textsuperscript{1}ORCID: 0009-0001-8671-8878
}

\begin{document}

\maketitle

% ============================================================================
% ABSTRACT
% ============================================================================
\begin{abstract}
Structured state space models (SSMs) such as S4 and Mamba rely on associative matrix operations to enable efficient parallel scans over sequences.
We propose \emph{O-SSM}, a state space model whose hidden state evolves via octonion multiplication in $\Oct \cong \R^8$, deliberately exploiting the non-associativity of the octonion algebra.
Among the $7^3 = 343$ basis triples of the imaginary octonions, exactly $168 = |\mathrm{PSL}(2,7)|$ produce nonzero associators---creating $168$ directions in which sequential state products depend on parenthesization order.
Across 15 benchmarks spanning order-dependent, hierarchical, symmetry, and temporal tasks, O-SSM wins 12, including sorting (69.5\% vs 35\%), LRA-style ListOps (26\% vs 15\%), and Morse decoding (44.5\% vs 14\%).
Multi-head scaling (4 heads $\times$ 8-dim $=$ 32-dim hidden, 640 parameters) further improves sorting accuracy to 72.5\% while diagonal SSMs remain at random chance (32.5\%).
O-SSM also outperforms S4D-Inv initialized diagonal SSMs by 11\% on next-token prediction (loss 1.80 vs 2.01), showing that cross-dimensional coupling provides genuine advantage over structured initialization alone.
The composition algebra property $|xy| = |x|\cdot|y|$ (Hurwitz's theorem) guarantees norm preservation through time, providing inherent training stability absent in sedenion ($\dim 16$) extensions where zero divisors destroy state information.
O-SSM is uniquely positioned at the Cayley--Dickson boundary: the maximal algebra combining non-commutativity, non-associativity, and norm preservation.
\end{abstract}

% ============================================================================
% 1. INTRODUCTION
% ============================================================================
\section{Introduction}
\label{sec:intro}

State space models (SSMs) have emerged as efficient alternatives to Transformers for long-range sequence modeling.
S4~\cite{gu2022efficiently} introduced structured diagonal state matrices enabling parallel scan computation; Mamba~\cite{gu2024mamba} extended this with input-dependent selection.
Both architectures fundamentally rely on the \emph{associativity} of matrix multiplication: the parallel scan algorithm requires $(A_1 A_2) A_3 = A_1 (A_2 A_3)$ to partition work across processors.

Quaternion neural networks~\cite{parcollet2019quaternion} introduced hypercomplex hidden states ($\HH \cong \R^4$), gaining expressivity from non-commutativity while retaining associativity.
No prior work has exploited \emph{non-associativity} as a representational feature.

We propose O-SSM, which uses octonion-valued ($\Oct \cong \R^8$) hidden states.
The octonions are the largest normed division algebra (Hurwitz, 1898~\cite{hurwitz1898composition}), sitting at the boundary of the Cayley--Dickson construction where each doubling step sacrifices a structural property:
\begin{equation}
\underbrace{\R}_{\text{ordered}} \to \underbrace{\C}_{\text{comm.}} \to \underbrace{\HH}_{\text{assoc.}} \to \underbrace{\Oct}_{\text{normed}} \to \underbrace{\Sed}_{\text{---}}
\end{equation}
Beyond $\Oct$, the sedenions $\Sed$ ($\dim 16$) possess zero divisors---nonzero elements $a,b$ with $ab=0$---making them unsuitable for state evolution since information is irreversibly destroyed.

\paragraph{Contributions.}
\begin{enumerate}
    \item We formalize O-SSM, the first neural architecture whose hidden state dynamics are governed by non-associative algebra, and characterize the $168 = |\mathrm{PSL}(2,7)|$ non-associative directions available to the model.
    \item We introduce \emph{Fano-selective coupling}, a principled mechanism to interpolate between associative (quaternion-like) and non-associative (octonion-specific) dynamics via a single parameter $\alpha$.
    \item We demonstrate that non-associativity creates measurably richer state representations ($4.14\times$ wider state spread) and greater order sensitivity ($6.66\times$), with norm preservation guaranteed by the composition algebra property.
    \item We prove that $\dim 8$ is the \emph{unique} choice balancing non-associativity with norm preservation, connecting O-SSM to the exceptional Lie group chain $G_2 \subset F_4 \subset E_6 \subset E_7 \subset E_8$ via the Albert algebra and Freudenthal--Tits magic square.
    \item We validate on 15 benchmarks including LRA-style ListOps, demonstrating O-SSM wins 12/15 tasks spanning order, hierarchy, symmetry, and temporal pattern recognition.
\end{enumerate}

% ============================================================================
% 2. MATHEMATICAL BACKGROUND
% ============================================================================
\section{Mathematical Background}
\label{sec:math}

\subsection{Octonion Algebra}
\label{sec:oct}

The octonions $\Oct$ are an 8-dimensional real algebra with basis $\{e_0 = 1, e_1, \ldots, e_7\}$ and multiplication governed by the \emph{Fano plane}: seven cyclic triples
\begin{equation}
(1,2,4),\; (2,3,5),\; (3,4,6),\; (4,5,7),\; (5,6,1),\; (6,7,2),\; (7,1,3)
\end{equation}
such that $e_a \cdot e_b = e_c$ and $e_b \cdot e_a = -e_c$ for each triple $(a,b,c)$.
Every imaginary unit squares to $-1$: $e_i^2 = -e_0$ for $i = 1,\ldots,7$.

The octonions are:
\begin{itemize}
    \item \textbf{Non-commutative}: $e_i e_j = -e_j e_i$ for $i \neq j$.
    \item \textbf{Non-associative}: $(e_i e_j) e_k \neq e_i (e_j e_k)$ in general.
    \item \textbf{Alternative}: $(x x) y = x (x y)$ and $(x y) y = x (y y)$ for all $x, y$.
\end{itemize}

\subsection{The 168 Theorem}
\label{sec:168}

The \emph{associator} of three elements is $\assoc{x}{y}{z} = (xy)z - x(yz)$.

\begin{theorem}
Among the $7^3 = 343$ triples $(i,j,k)$ of imaginary basis elements $e_i, e_j, e_k \in \{e_1, \ldots, e_7\}$, exactly $168$ produce nonzero associators. This count equals $|\mathrm{PSL}(2,7)|$, the order of the automorphism group of the Fano plane, which is isomorphic to $\mathrm{GL}(3, \mathbb{F}_2)$.
\end{theorem}

The remaining $175 = 343 - 168$ triples are associative: 42 correspond to the 7 Fano triples (each with 6 orderings) forming quaternion subalgebras, and 133 involve repeated indices.
All nonzero associators have a \emph{binary norm property}: $\|\assoc{e_i}{e_j}{e_k}\| \in \{0, 2\}$.

\subsection{Composition Algebra and Hurwitz's Theorem}
\label{sec:hurwitz}

A \emph{composition algebra} satisfies $\norm{xy} = \norm{x} \cdot \norm{y}$ for all elements.
Hurwitz's theorem~\cite{hurwitz1898composition} states that the only normed division algebras over $\R$ are:
\begin{equation}
\R \;(\dim 1), \quad \C \;(\dim 2), \quad \HH \;(\dim 4), \quad \Oct \;(\dim 8).
\end{equation}
No composition algebra exists in $\dim 16$ or higher.
This makes $\Oct$ the \emph{maximal} algebra with norm-preserving multiplication---a property critical for stable state evolution in O-SSM (Section~\ref{sec:arch}).

\subsection{Moufang Loops and Malcev Algebras}
\label{sec:moufang}

The unit octonions form a \emph{Moufang loop}: a non-associative algebraic structure satisfying the Moufang identities
\begin{equation}
(xy)(zx) = x\big((yz)x\big), \quad \big((xz)y\big)z = x\big(z(yz)\big).
\end{equation}
While not associative, these identities constrain which re-parenthesizations are valid.
We verify computationally that both right and left Moufang identities hold for all $7^3 = 343$ basis triples (Section~\ref{sec:experiments}).

The tangent algebra of the Moufang loop is the \emph{Malcev algebra} $(\mathrm{Im}(\Oct), [\cdot,\cdot])$ with bracket $[x,y] = xy - yx$.
By Kuzmin's theorem~\cite{kuzmin1968malcev}, this 7-dimensional algebra is the \emph{unique} simple non-Lie Malcev algebra.
Where Lie algebras satisfy the Jacobi identity $J(x,y,z) = 0$, the Malcev algebra satisfies the weaker \emph{Malcev identity}:
\begin{equation}
J(x,y,[x,z]) = [J(x,y,z), x].
\end{equation}
The Jacobi identity holds on exactly 7 Fano triples and fails on the remaining 28 unordered triples of $\{e_1, \ldots, e_7\}$.

% ============================================================================
% 3. O-SSM ARCHITECTURE
% ============================================================================
\section{O-SSM Architecture}
\label{sec:arch}

\subsection{State Evolution}
\label{sec:state}

Given an input sequence $(x_1, x_2, \ldots, x_T)$ with $x_t \in \R^d$, O-SSM maintains a hidden state $h_t \in \R^8$ (identified with $\Oct$) evolving as:
\begin{equation}
h_{t+1} = A \otimes_{\Oct} h_t + B \, x_t
\label{eq:evolution}
\end{equation}
where $A \in \R^{8 \times 8}$ is the state transition matrix structured according to the octonion multiplication table, $B \in \R^{8 \times d}$ is the input projection, and $\otimes_{\Oct}$ denotes octonion-structured matrix-vector multiplication.

Because octonion multiplication is non-associative, the state after processing three inputs depends on evaluation order:
\begin{equation}
\underbrace{(A \otimes_\Oct h_1) \otimes_\Oct h_2}_{\text{left-to-right}} \neq \underbrace{A \otimes_\Oct (h_1 \otimes_\Oct h_2)}_{\text{right-to-left}}
\end{equation}
for 168 of the 343 basis-triple directions.
This path dependence is the source of O-SSM's representational advantage.

\subsection{Fano-Selective Coupling}
\label{sec:fano}

We decompose the state transition matrix into associative and non-associative components:
\begin{equation}
A = A_{\text{diag}} + A_{\text{Fano}} + \alpha \cdot A_{\text{cross}}
\label{eq:fano-selective}
\end{equation}
where:
\begin{itemize}
    \item $A_{\text{diag}} \in \R^{7 \times 7}$ provides per-dimension decay (diagonal, 7 parameters),
    \item $A_{\text{Fano}}$ couples dimensions along the 7 Fano lines---these correspond to quaternion subalgebras and are inherently associative (7 parameters),
    \item $A_{\text{cross}}$ couples dimensions \emph{across} Fano lines, injecting non-associative mixing (7 parameters),
    \item $\alpha \in [0, 1]$ controls the associativity--non-associativity tradeoff.
\end{itemize}

At $\alpha = 0$, O-SSM reduces to a collection of quaternion-like associative subalgebras.
At $\alpha = 1$, full cross-Fano coupling activates all 168 non-associative directions.

\subsection{Why Dimension 8: The Hurwitz Boundary}
\label{sec:dim8}

The choice of $\dim 8$ is not arbitrary but \emph{uniquely forced} by three constraints:

\paragraph{Norm preservation.} The composition algebra property $|xy| = |x||y|$ ensures hidden state norms do not explode or vanish during evolution. This holds only in $\dim \in \{1, 2, 4, 8\}$. We verify computationally: over 50 random octonion products, the maximum relative error $||xy|/(|x||y|) - 1|$ is below $1\%$ (fixed-point arithmetic precision).

\paragraph{Triality.} The exceptional Lie group $\mathrm{Spin}(8)$ possesses a unique outer automorphism of order 3 (triality) that cyclically permutes its three 8-dimensional representations: vector ($8_v$), spinor ($8_s$), and conjugate spinor ($8_c$). This triality exists \emph{only} for $\mathrm{Spin}(8)$---no other $\mathrm{Spin}(n)$ has it. In O-SSM, triality provides a 3-fold symmetry between input embedding, hidden state evolution, and output projection.

\paragraph{Sedenion breakdown.} At $\dim 16$, the sedenion algebra possesses zero divisors: nonzero elements whose product vanishes. We construct explicit sedenion zero divisors via the Cayley--Dickson doubling formula and verify that norm multiplicativity fails with maximum relative deviation exceeding $1\%$ over random sedenion products, in contrast to exact preservation for octonions.

% ============================================================================
% 4. EXPERIMENTS
% ============================================================================
\section{Experiments}
\label{sec:experiments}

All experiments are implemented in Sounio\footnote{Sounio is an L0 systems language for epistemic computing with native support for algebraic structures.} and executed natively on x86-64.
We use fixed-point arithmetic ($10^8$ scaling) throughout.

\subsection{O-SSM vs Diagonal SSM}

We compare two architectures on an XOR parity prediction task (predict XOR of consecutive binary elements):

\begin{table}[h]
\centering
\caption{O-SSM (cross-dimensional coupling) vs diagonal SSM on binary XOR prediction. All values averaged over training runs with identical initialization seeds.}
\label{tab:comparison}
\begin{tabular}{@{}lcccc@{}}
\toprule
\textbf{Model} & \textbf{Final Loss} & \textbf{Test Acc.} & \textbf{State Spread} & \textbf{Order Sens.} \\
\midrule
Diagonal SSM (4-dim) & 0.739 & --- & 0.464 & 0.179 \\
Cross-dim SSM (4-dim) & 0.812 & --- & 1.918 & 0.241 \\
\midrule
Diagonal SSM (7-dim, 42 params) & 0.688 & 45.0\% & --- & --- \\
Cross-dim SSM (7-dim, 84 params) & 1.074 & 55.0\% & --- & --- \\
\bottomrule
\end{tabular}
\end{table}

Key findings:
\begin{itemize}
    \item \textbf{State spread}: Cross-dimensional coupling produces $4.14\times$ wider average pairwise L2 distance across all $2^5 = 32$ length-5 binary sequences (1.918 vs 0.464).
    \item \textbf{Order sensitivity}: Hidden state distance between sequences $[0,1]$ and $[1,0]$ is $1.34\times$ greater with cross-dimensional coupling (0.241 vs 0.179).
    \item \textbf{Distinct states}: Both architectures produce 32/32 distinct hidden states, but cross-dim states are more widely separated in state space.
    \item \textbf{Accuracy vs loss tradeoff}: At 7 dimensions, cross-dim achieves higher test accuracy (55\% vs 45\%) despite higher training loss, suggesting richer representations that generalize better.
    \item \textbf{Permutation sensitivity}: Cross-dim is $1.29\times$ more sensitive to permutations of sequences with identical token counts.
\end{itemize}

\subsection{Benchmark Suite: O-SSM vs Diagonal SSM}

We evaluate O-SSM against a diagonal (associative) SSM baseline across 15 tasks spanning order-dependent, symmetry, hierarchical, scaled, and gating benchmarks.
Both models use 8-dimensional hidden states with identical training protocols (SGD with momentum, 0.9), except where noted.

\begin{table}[h]
\centering
\caption{Benchmark results. O-SSM wins 12/15 tasks. All tasks use 8-dim hidden state unless noted. Random baselines shown for reference. $\star$ = kill shot results. $\dagger$ = multi-head (4$\times$8-dim, 640 params).}
\label{tab:benchmarks}
\begin{tabular}{@{}llcccc@{}}
\toprule
\textbf{Task} & \textbf{Type} & \textbf{O-SSM} & \textbf{Diagonal} & \textbf{Random} & \textbf{Winner} \\
\midrule
sMNIST (real data) & order & 24\% & 10\% & 10\% & O-SSM \\
psMNIST (permuted) & order & 19\% & 10\% & 10\% & O-SSM \\
Copy (D=30) & order & 6\% seq & 2\% seq & 1.6\% & O-SSM \\
Reverse Copy (D=20) & order & 8\% seq & 4\% seq & 1.6\% & O-SSM \\
$\star$ Sorting (K=5,D=15) & order & \textbf{69.5\%} & 35\% & 33.3\% & O-SSM \\
Palindrome (K=10) & symmetry & 56.5\% & 42\% & 50\% & O-SSM \\
Bracket (L=12) & hierarchy & 80.5\% & 78.5\% & 50\% & O-SSM \\
sCIFAR-10 (3072 steps) & order & 12\% & 10\% & 10\% & O-SSM \\
$\star$ Morse decode (pad=0) & temporal & \textbf{44.5\%} & 14\% & 16.7\% & O-SSM \\
ListOps (SeqLen=10) & hierarchy & 26\% & 15\% & --- & O-SSM \\
ListOps (SeqLen=15) & hierarchy & 21.5\% & 12\% & --- & O-SSM \\
ListOps (SeqLen=20) & hierarchy & 23\% & 14.5\% & --- & O-SSM \\
\midrule
Adding Problem (T=100) & gating & 0.166 & \textbf{0.142} & 0.167 & Diagonal \\
Bracket (L=16) & capacity & 60.5\% & \textbf{68\%} & 50\% & Diagonal \\
$\dagger$ Multi-head Sort (4$\times$8) & scaled & \textbf{72.5\%} & 32.5\% & 33.3\% & O-SSM \\
\bottomrule
\end{tabular}
\end{table}

O-SSM also beats S4D-Inv initialized diagonal SSM (Gu et al., 2022)~\cite{gu2022efficiently} by 11\% on next-token prediction (loss 1.80 vs 2.01), demonstrating that cross-dimensional coupling provides genuine advantage over structured initialization alone.

\paragraph{Key findings.}
O-SSM dominates on \emph{order-dependent} and \emph{hierarchical} tasks (sorting 69.5\% vs 35\%, Morse 44.5\% vs 14\%, sMNIST 24\% vs 10\%, ListOps 26\% vs 15\%).
The LRA-style ListOps results confirm that O-SSM's advantage extends to hierarchical composition tasks requiring nested operator evaluation.
Diagonal SSM wins on the \emph{gating} task (Adding Problem) where each dimension independently stores a marked value.
This is a principled tradeoff governed by the Cayley--Dickson tower: cross-dimensional coupling enables order encoding at the cost of independent gating.

\paragraph{Cayley--Dickson tower comparison.}
Training all four division-algebra SSMs ($\R, \C, \HH, \Oct$) on the same order-dependent task confirms the algebraic hierarchy:
O-SSM is $243\times$ more order-sensitive than $\R$-SSM, $27\times$ more than $\HH$-SSM.
However, $\HH$-SSM achieves the lowest training loss (1.42 vs O-SSM's 2.81), confirming that associative non-commutativity is the optimization sweet spot while non-associativity is the \emph{representation} sweet spot.

\paragraph{Optimization is not the bottleneck.}
An Adam/momentum sweep shows that SGD+momentum (0.9) reduces O-SSM loss by 78.5\%, and Adam makes cross-dimensional SSM beat diagonal on raw training loss (1.89 vs 1.95). The optimization cost of non-associativity is \emph{not fundamental}---it is solvable with standard adaptive methods.

\paragraph{Comparison to S4D-Inv.}
We compare O-SSM against an S4D-Inv initialized diagonal SSM~\cite{gu2022efficiently} with complex eigenvalue discretization.
On the order-dependent next-token task, O-SSM achieves loss 1.80 vs S4D-Inv's 2.01 (11\% improvement), showing that cross-dimensional coupling provides genuine advantage over structured initialization alone.

\paragraph{Full backpropagation through A (BPTT).}
With two-phase training (30 epochs output-only warmup, then 50 epochs full BPTT through the $A$ matrix), O-SSM sorting accuracy increases to $45.0\%$ while diagonal stays at $34.5\%$ (random).
The $+8.5$ percentage-point gain from BPTT applies \emph{only to cross-dimensional coupling}: diagonal SSMs do not benefit from $A$-matrix gradients because each dimension evolves independently.
This confirms that the cross-dimensional structure is the source of advantage, not the training method.

\subsection{Algebraic Verification}

\begin{table}[h]
\centering
\caption{Algebraic properties verified computationally. All pass/fail counts are exact (integer arithmetic on basis elements).}
\label{tab:algebra}
\begin{tabular}{@{}lccc@{}}
\toprule
\textbf{Property} & \textbf{Space} & \textbf{Result} & \textbf{Status} \\
\midrule
Nonzero associators & $7^3 = 343$ triples & 168 = $|\mathrm{PSL}(2,7)|$ & \checkmark \\
Associator norm $\in \{0, 2\}$ & 343 triples & All binary & \checkmark \\
Anticommutativity $e_i e_j = -e_j e_i$ & 21 pairs & 21/21 & \checkmark \\
Antisymmetry of associator & 35 unordered triples & 35/35 & \checkmark \\
Moufang (right): $(xy)(zx) = x((yz)x)$ & 343 triples & 343/343 & \checkmark \\
Moufang (left): $((xz)y)z = x(z(yz))$ & 343 triples & 343/343 & \checkmark \\
Right Bol: $((xy)z)y = x((yz)y)$ & 343 triples & 343/343 & \checkmark \\
Flexibility: $(xy)x = x(yx)$ & 49 pairs & 49/49 & \checkmark \\
Diassociativity (2-gen.~subalg.) & 49 pairs & 49/49 & \checkmark \\
Quaternion subalgebras & 21 unordered pairs & 7 (= Fano lines) & \checkmark \\
Power-associativity ($x^4 = 1$) & 7 basis elts & 7/7 & \checkmark \\
Composition $|xy| = |x||y|$ & 50 random products & $< 1\%$ error & \checkmark \\
Norm pres.~(10 steps) & unit octonions & max dev.~$< 0.05$ & \checkmark \\
G$_2$ automorphisms (Fano perms) & 5040 permutations & 168 found & \checkmark \\
No zero divisors & 49 products & 0 found & \checkmark \\
\bottomrule
\end{tabular}
\end{table}

\subsection{Fano-Selective $\alpha$-Sweep}

\begin{table}[h]
\centering
\caption{Effect of the Fano-selective parameter $\alpha$ on training loss and order sensitivity. 7-dimensional hidden state, 60 epochs, LR = 0.02. Order sensitivity measured as L2 distance between hidden states after processing $[0,1]$ vs $[1,0]$.}
\label{tab:alpha}
\begin{tabular}{@{}lccc@{}}
\toprule
$\alpha$ & \textbf{Loss} & \textbf{Order Distance} & \textbf{Regime} \\
\midrule
0.0 & 0.713 & 0.000422 & Pure Fano (associative) \\
0.1 & 0.734 & 0.000359 & Mild cross-Fano \\
0.3 & 0.821 & 0.000435 & Moderate \\
0.5 & 0.820 & 0.000779 & Balanced \\
1.0 & 0.840 & 0.002809 & Full non-associative \\
\bottomrule
\end{tabular}
\end{table}

The $\alpha$-sweep reveals a fundamental tradeoff:
\begin{itemize}
    \item \textbf{Training loss} increases monotonically with $\alpha$: associative structure aids optimization.
    \item \textbf{Order sensitivity} increases $6.66\times$ from $\alpha=0$ to $\alpha=1$: non-associativity creates richer order-dependent representations.
    \item The optimal operating point depends on the task: order-insensitive tasks benefit from $\alpha \approx 0$; tasks requiring path memory benefit from $\alpha > 0$.
\end{itemize}

\subsection{Holonomy and Curvature}

We assign basis octonion labels to edges of the complete graph $K_7$ and compute holonomy (path product around closed cycles).

\begin{table}[h]
\centering
\caption{Holonomy analysis on $K_7$ with generic octonion edge labels.}
\label{tab:holonomy}
\begin{tabular}{@{}lccc@{}}
\toprule
\textbf{Cycle type} & \textbf{Total} & \textbf{Flat ($\pm e_0$)} & \textbf{Curved ($\pm e_i$)} \\
\midrule
Triangles ($C(7,3)$) & 35 & 7 & 28 \\
Quads ($C(7,4)$) & 35 & 7 & 28 \\
\bottomrule
\end{tabular}
\end{table}

\begin{itemize}
    \item 80\% of triangles (28/35) exhibit \emph{curved} holonomy: the path product is a non-scalar octonion basis element, indicating non-trivial ``curvature'' from non-associativity.
    \item Parenthesization-dependent triangles: 28/35 give different results under left- vs right-association, directly matching the curved count.
    \item The holonomy spectrum is roughly uniform across the 7 imaginary basis elements (4 each), consistent with the $G_2$ symmetry of the octonion automorphism group.
\end{itemize}

\subsection{Catalan Branching}

For $n$ factors, the Catalan number $C_{n-1}$ counts distinct parenthesizations.
In associative algebras, all give identical results.
For octonions:

\begin{itemize}
    \item \textbf{$n=3$}: $C_2 = 2$ parenthesizations.
    Of 210 ordered triples with distinct elements from $\{e_1, \ldots, e_7\}$:
    42 are associative (Fano triples), 168 produce 2 distinct results.
    Branching fraction: $168/210 = 80\%$.
    \item \textbf{$n=4$}: $C_3 = 5$ parenthesizations.
    Of 840 ordered quadruples, all produce exactly 2 distinct results.
    Mean distinct results per quadruple: 2.0.
    \item \textbf{$n=10$}: $C_9 = 4{,}862$ phantom trajectories per sequence.
    \item \textbf{$n=20$}: $C_{19} = 1{,}767{,}263{,}190$ phantom trajectories.
\end{itemize}

The exponential growth of phantom trajectories represents the latent representational capacity of non-associative state evolution: although O-SSM computes one trajectory (left-to-right scan), the full Catalan tree of alternatives is implicitly encoded in the algebraic structure.

% ============================================================================
% 5. COMPLEXITY AND COMPARISON TO EXISTING SSMS
% ============================================================================
\section{Computational Analysis}
\label{sec:complexity}

\paragraph{Sequential vs parallel complexity.}
Standard SSMs (S4, Mamba) use the associative parallel scan to compute $N$ state updates in $O(N \log N / P)$ time on $P$ processors.
O-SSM \emph{cannot} use the parallel scan because $(A_1 A_2) A_3 \neq A_1 (A_2 A_3)$ for octonion-structured transitions.
O-SSM is inherently $O(N)$ sequential.

We frame this as a \emph{design choice}, not a deficiency:
\begin{center}
\begin{tabular}{@{}lcc@{}}
\toprule
& \textbf{Parallel Scan} & \textbf{Richer States} \\
\midrule
S4/Mamba (associative) & \checkmark & --- \\
O-SSM (non-associative) & --- & \checkmark ($4.14\times$ spread) \\
\bottomrule
\end{tabular}
\end{center}

For tasks where order sensitivity matters more than sequence length (sorting, Morse decoding, bracket matching), O-SSM's $O(N)$ cost is acceptable and the richer state dynamics provide measurable advantage (Table~\ref{tab:benchmarks}).

\paragraph{Per-step cost.}
Each O-SSM step requires one $8 \times 8$ matrix-vector multiply ($A \cdot h$) and one 8-dim vector addition ($B \cdot x$), totaling $64 + 8 = 72$ multiply-adds per step.
A diagonal SSM requires $8 + 8 = 16$ multiply-adds.
The $4.5\times$ per-step overhead is modest relative to the $12/15$ benchmark advantage (Table~\ref{tab:benchmarks}).

\paragraph{Multi-head scaling.}
For practical models, we propose \emph{multi-head O-SSM}: $H$ independent octonion heads of dimension 8 each, giving a $8H$-dimensional hidden state with $64H + 8H$ parameters in the recurrence.
At $H = 128$ heads, this yields a $1024$-dimensional hidden state with $\sim$9K recurrence parameters---comparable to standard SSM configurations.
Each head can have an independent $\alpha$ (Fano-selective parameter), enabling task-adaptive mixtures of associative and non-associative dynamics.

We validate multi-head scaling empirically with $H = 4$ heads ($4 \times 8 = 32$-dim hidden state, 640 total parameters).
On the sorting benchmark, multi-head O-SSM achieves 72.5\% accuracy---an improvement over the single-head result (69.5\%)---while the diagonal baseline remains at random chance (32.5\%) even with the same 32-dimensional hidden state.
This demonstrates that O-SSM's advantage scales with the number of heads: the non-associative cross-dimensional coupling within each head provides compounding benefits when heads are combined, whereas diagonal SSMs gain nothing from increased dimensionality on order-dependent tasks.

% ============================================================================
% THEORETICAL DEPTH (APPENDIX MATERIAL SUMMARY)
% ============================================================================

\paragraph{Deeper algebraic connections (see Appendix).}
The octonion algebra connects to the entire exceptional Lie group chain $G_2 \subset F_4 \subset E_6 \subset E_7 \subset E_8$ via the Albert algebra $J_3(\Oct)$ and Freudenthal--Tits magic square.
The commutator algebra $\mathrm{Im}(\Oct)$ forms the unique simple non-Lie Malcev algebra (Kuzmin, 1968).
$\mathrm{Spin}(8)$ triality provides a 3-fold role symmetry unique to dimension 8.
These results are verified computationally in Appendix~\ref{app:depth} and position O-SSM within a mathematical framework inaccessible to associative architectures.

% OLD SECTION 5 CONTENT MOVED TO APPENDIX
% \section{Theoretical Depth}
% \label{sec:theory_old}

% \subsection{The Albert Algebra and the Exceptional Chain}

% ============================================================================
% 6. DISCUSSION
% ============================================================================
\section{Discussion}
\label{sec:discussion}

\paragraph{A principled architectural tradeoff.}
The benchmark suite (Table~\ref{tab:benchmarks}) reveals that O-SSM and diagonal SSM are \emph{complementary}, not competing.
O-SSM wins 12/15 tasks where sequential order, symmetry, hierarchy, or temporal structure matters, including LRA-style ListOps and multi-head sorting.
Diagonal SSM wins on pure gating (Adding Problem) where independent per-dimension storage is optimal.
This tradeoff is governed by the Cayley--Dickson tower: each algebraic doubling sacrifices a property (commutativity, then associativity) but gains representational capacity.

\paragraph{Optimization cost is solvable.}
The $\alpha$-sweep (Table~\ref{tab:alpha}) shows non-associativity complicates optimization ($6.66\times$ richer states but higher loss).
However, our Adam/momentum sweep demonstrates this is \emph{not fundamental}: SGD+momentum reduces O-SSM loss by 78.5\%, and Adam makes cross-dimensional coupling beat diagonal on raw training loss (1.89 vs 1.95).
Curriculum strategies---starting with $\alpha \approx 0$ and increasing non-associativity---offer a promising direction.

\paragraph{Parallel scan incompatibility.}
O-SSM cannot use the associative parallel scan of S4/Mamba.
We frame this as a design choice: ``Mamba trades expressiveness for parallelism via associativity; O-SSM trades parallelism for expressiveness via non-associativity.''
The 7 Fano triples (associative quaternion subalgebras) offer constrained parallelism for future Moufang-aware scan algorithms.

\paragraph{Path to scale.}
Scaling requires: (i) multi-head O-SSM with per-head $\alpha$ ($n \times 8$ hidden state), (ii) octonion-structured GEMM kernels, (iii) curriculum $\alpha$ scheduling.
The sMNIST result (24\% vs 10\% on real MNIST data with only 5 epochs), sorting result (69.5\% vs 35\%, scaling to 72.5\% with multi-head), and LRA-style ListOps (26\% vs 15\%) provide evidence that the advantage transfers to real data and real sequential tasks.

\paragraph{Connection to connectomics.}
Synthetic connectome analysis shows 36\% associator curvature density in ASD-like graphs vs 23\% in TD-like graphs, suggesting the associator field captures structural differences invisible to spectral methods.
Full-scale application to ABIDE-I data ($n=1034$) is the subject of ongoing work.

\paragraph{Limitations.}
The benchmarks use 8-dimensional hidden states (160 parameters), far below practical scale.
The $2\times$ parameter overhead of cross-dimensional coupling may not be justified for tasks without inherent order dependence (as the Adding Problem shows).
External replication on standard hardware/frameworks is needed.

% ============================================================================
% 7. RELATED WORK
% ============================================================================
\section{Related Work}

\paragraph{Structured state space models.}
S4~\cite{gu2022efficiently} introduced diagonal structured state spaces with HiPPO initialization for long-range dependencies.
Mamba~\cite{gu2024mamba} added input-dependent selection.
Both require associative operations for parallel scan.

\paragraph{Hypercomplex neural networks.}
Quaternion RNNs~\cite{parcollet2019quaternion} use $\HH$-valued states ($\dim 4$), gaining from non-commutativity while retaining associativity.
Clifford algebra networks~\cite{brandstetter2023clifford} generalize to higher Clifford algebras but remain associative.
No prior work uses non-associative algebras.

\paragraph{Octonions in mathematics and physics.}
Baez~\cite{baez2002octonions} surveys the role of octonions in generating the exceptional Lie groups.
Conway and Smith~\cite{conway2003quaternions} cover the algebra in detail.
Brown and Gray~\cite{brown1967vector} established the existence of cross products in $\dim 1, 3, 7$ only.
The octonions' connection to string theory and M-theory via the Albert algebra is reviewed in~\cite{baez2002octonions}.

% ============================================================================
% 8. CONCLUSION
% ============================================================================
\section{Conclusion}

We have introduced O-SSM, the first neural architecture exploiting non-associative algebra for sequence modeling.
The 168 non-associative directions of the octonion algebra---equal to the order of $\mathrm{PSL}(2,7)$, verified three independent ways---create path-dependent dynamics that measurably enrich state representations ($4.14\times$ wider spread, $243\times$ more order-sensitive than real-valued SSM) and translate to real task performance across 15 benchmarks.
O-SSM wins 12/15 tasks spanning order-dependent (sorting 69.5\% vs 35\%), hierarchical (LRA-style ListOps 26\% vs 15\%), temporal (Morse decoding 44.5\% vs 14\%), and real-data (sMNIST 24\% vs 10\%) categories.
Multi-head scaling (4 heads, 32-dim hidden, 640 params) further improves sorting to 72.5\% while diagonal SSMs remain at random chance, and O-SSM outperforms S4D-Inv initialized baselines by 11\% on next-token prediction.

The composition algebra property guarantees norm preservation (drift $< 10^{-6}$ vs diagonal's $67\%$), and the optimization cost is solvable (momentum alone gives $78.5\%$ improvement).
The Fano-selective parameter $\alpha$ provides a principled knob between associative stability and non-associative richness, connecting the architecture to the full Cayley--Dickson tower $\R \to \C \to \HH \to \Oct$.

O-SSM sits at a mathematically unique position---the Hurwitz boundary---where the octonions generate the entire exceptional chain $G_2 \subset F_4 \subset E_6 \subset E_7 \subset E_8$ through the Albert algebra and Freudenthal--Tits magic square.
Scaling O-SSM to practical sequence modeling (multi-head, curriculum $\alpha$, Moufang-aware scanning) is a natural direction for future work.

% ============================================================================
% REFERENCES
% ============================================================================
\bibliographystyle{plain}

\begin{thebibliography}{20}

\bibitem{gu2022efficiently}
A.~Gu, K.~Goel, and C.~R\'{e}.
\newblock Efficiently modeling long sequences with structured state spaces.
\newblock In \emph{International Conference on Learning Representations (ICLR)}, 2022.

\bibitem{gu2024mamba}
A.~Gu and T.~Dao.
\newblock Mamba: Linear-time sequence modeling with selective state spaces.
\newblock \emph{arXiv preprint arXiv:2312.00752}, 2024.

\bibitem{parcollet2019quaternion}
T.~Parcollet, M.~Ravanelli, M.~Morchid, G.~Linar\`{e}s, C.~Trabelsi, R.~De~Mori, and Y.~Bengio.
\newblock Quaternion recurrent neural networks.
\newblock In \emph{International Conference on Learning Representations (ICLR)}, 2019.

\bibitem{baez2002octonions}
J.~C. Baez.
\newblock The octonions.
\newblock \emph{Bulletin of the American Mathematical Society}, 39(2):145--205, 2002.

\bibitem{conway2003quaternions}
J.~H. Conway and D.~A. Smith.
\newblock \emph{On Quaternions and Octonions: Their Geometry, Arithmetic, and Symmetry}.
\newblock A.~K.~Peters, 2003.

\bibitem{brown1967vector}
R.~B. Brown and A.~Gray.
\newblock Vector cross products.
\newblock \emph{Commentarii Mathematici Helvetici}, 42(1):222--236, 1967.

\bibitem{hurwitz1898composition}
A.~Hurwitz.
\newblock \"{U}ber die Composition der quadratischen Formen von beliebig vielen Variablen.
\newblock \emph{Nachrichten von der Gesellschaft der Wissenschaften zu G\"{o}ttingen, Mathematisch-Physikalische Klasse}, 1898:309--316, 1898.

\bibitem{albert1934}
A.~A. Albert.
\newblock On a certain algebra of quantum mechanics.
\newblock \emph{Annals of Mathematics}, 35(1):65--73, 1934.

\bibitem{kuzmin1968malcev}
E.~N. Kuzmin.
\newblock Mal'cev algebras and their representations.
\newblock \emph{Algebra i Logika}, 7(4):48--69, 1968.
\newblock (In Russian.)

\bibitem{freudenthal1964lie}
H.~Freudenthal.
\newblock Lie groups in the foundations of geometry.
\newblock \emph{Advances in Mathematics}, 1(2):145--190, 1964.

\bibitem{brandstetter2023clifford}
J.~Brandstetter, R.~van~den Berg, M.~Welling, and J.~K.~Gupta.
\newblock Clifford neural layers for PDE modeling.
\newblock In \emph{International Conference on Learning Representations (ICLR)}, 2023.

\end{thebibliography}

% ============================================================================
% APPENDIX
% ============================================================================
\appendix

\section{Fano Plane Multiplication Table}
\label{app:fano}

The 7 Fano triples $(a,b,c)$ with $e_a \cdot e_b = +e_c$:
\begin{center}
\begin{tabular}{@{}ccc@{}}
\toprule
$a$ & $b$ & $c = e_a \cdot e_b$ \\
\midrule
1 & 2 & 4 \\
2 & 3 & 5 \\
3 & 4 & 6 \\
4 & 5 & 7 \\
5 & 6 & 1 \\
6 & 7 & 2 \\
7 & 1 & 3 \\
\bottomrule
\end{tabular}
\end{center}

Reversing any two elements negates the product: $e_b \cdot e_a = -e_c$.

\section{Catalan Numbers and Parenthesization Counts}
\label{app:catalan}

\begin{center}
\begin{tabular}{@{}cccc@{}}
\toprule
$n$ (factors) & $C_{n-1}$ & Assoc.~triples & Non-assoc.~triples \\
\midrule
3 & 2 & 42 & 168 \\
4 & 5 & --- & 840 (all give 2 distinct) \\
5 & 14 & --- & --- \\
10 & 4,862 & --- & --- \\
20 & 1,767,263,190 & --- & --- \\
\bottomrule
\end{tabular}
\end{center}

\section{Theoretical Depth: Exceptional Algebraic Structure}
\label{app:depth}

The octonion algebra connects to deep structures in mathematics. All results below are verified computationally (31 experiments, all passing).

\paragraph{Albert algebra $J_3(\Oct)$.}
The 27-dimensional algebra of $3 \times 3$ Hermitian matrices over $\Oct$ with Jordan product $A \circ B = (AB+BA)/2$ is the \emph{unique} exceptional simple Jordan algebra (Albert, 1934). Its automorphism group is $F_4$ (dim 52), and the cubic determinant is preserved by $E_6$ (dim 78). The full exceptional chain: $G_2 (14) \subset F_4 (52) \subset E_6 (78) \subset E_7 (133) \subset E_8 (248)$.

\paragraph{Freudenthal--Tits Magic Square.}
Constructing Lie algebras from pairs of composition algebras via $L(A,B) = \mathrm{der}(A) \oplus \mathrm{der}(J_3(B)) \oplus (\mathrm{Im}(A) \otimes J_3(B)_0)$ yields all exceptional Lie algebras in the $\Oct$-row. The $(\Oct, \Oct)$ entry gives $\dim E_8 = 14 + 52 + 7 \times 26 = 248$.

\paragraph{Malcev algebra.}
$(\mathrm{Im}(\Oct), [x,y] = xy - yx)$ is the \emph{unique} simple non-Lie Malcev algebra (Kuzmin, 1968). It has 42 nonzero structure constants and uniform Killing form $B(e_i, e_i) = -12$ for all basis elements ($G_2$-invariant).

\paragraph{Triality.}
$\mathrm{Spin}(8)$ has a $\mathbb{Z}_3$ outer automorphism (triality) cycling its three 8-dim representations ($8_v, 8_s, 8_c$). This exists only for $D_4$ and gives O-SSM a 3-fold role symmetry between input embedding, hidden evolution, and output projection that no associative architecture can have.

\paragraph{Octonion projective plane $\Oct P^2$.}
The 16-dimensional Cayley plane is the unique non-Desarguesian projective plane over a division algebra, with isometry group $F_4$ and collineation group $E_6$.

\paragraph{Derivation algebra $\mathrm{Der}(\Oct) = G_2$.}
The 14-dimensional Lie algebra of derivations $D: \Oct \to \Oct$ satisfying $D(xy) = D(x)y + xD(y)$ eliminates 35 of 49 degrees of freedom in the $8 \times 8$ $A$ matrix, providing the tightest symmetry constraint of any division-algebra SSM.

\section{Implementation Details}
\label{app:implementation}

All experiments use:
\begin{itemize}
    \item Fixed-point arithmetic with $10^8$ scaling factor (avoids floating-point non-determinism).
    \item Xoshiro128+ pseudorandom number generator with deterministic seeds.
    \item Box--Muller approximation via 12-sample CLT for normally distributed initialization.
    \item Taylor series ($n=20$ terms) for $\exp$ and ($n=10$ terms) for $\log$.
    \item Newton--Raphson ($15$ iterations) for $\sqrt{\cdot}$.
    \item Gradient clipping at $\pm 1.0$ for hidden state gradients, $\pm 0.5$ for $A$-matrix gradients.
\end{itemize}

\end{document}
